\section{Funciones}

Similar a como las variables son nombres que se asignan a valores, las funciones son nombres que se asignan a procedimientos. Esto permite pensar y escribir un pedazo de código sola una vez y luego reutilizar esa misma lógica múltiples veces.

Distinguimos entre dos tipos de funciones: las funciones \emph{\textit{built-in}} (integradas o nativas), que son funciones que creó alguien más y ya vienen predefinidas en Python, y las funciones \emph{\textit{user-defined}} (definidas por el usuario), que son aquellas que creamos nosotros.

\subsection{Funciones de conversión de tipos}

% TODO: Y type

\subsection{Algunas funciones matemáticas}

% TODO: max, min, abs, pow, round

\subsection{Funciones de entrada y salida}

Las funciones de entrada y salida se utilizan en el programa para interactuar con el usuario.
% TODO: En la consola!! Introducir consola en algún punto. Probablemente ya debería haberlo hecho hace rato. 

\subsubsection{Función \mintinline{python}|print|}

La función \mintinline{python}|print| (que traduce a imprima) recibe un argumento y lo muestra en la consola.

\subsubsection{Función \mintinline{python}|input|}

La función \mintinline{python}|input| se utiliza para leer un valor de la consola. Es decir, se utiliza para que el usuario ingrese un valor y se lo asigne a una variable.

\subsection{}

% TODO: ord, chr, help